\chapter{}
Entre os clientes a quem Paulo dava assistência na implantação e manutenção do cultivo de seringueiras, estava a Companhia de Seguros Sul América.
Esse grupo andava negociando com o governo do Piauí um valioso contrato mediante o qual segurariam uma grande parte dos bens do estado, em troca do compromisso de implantar na região de Miguel Alves, cidadezinha localizada à beira do Rio Parnaíba, um moderno complexo agropecuário compreendendo o plantio de 2500 hectares de arroz irrigado e uma usina de beneficiamento e empacotamento do produto que empregasse a mais avançada tecnologia disponível no mercado nacional.
A Sul América era mais uma das empresas que, seduzidas pelas facilidades oferecidas por órgãos de fomento do desenvolvimento regional como SUDAM e SUDENE, passara a investir no agronegócio, adquirindo fazendas em Minas Gerais, Pará e agora, no Piauí.
Para atender ao novo ramo de atividades foi necessário formar novos quadros e foi assim que chegou até Paulo, cujo trabalho eles tiveram oportunidade de conhecer por causa da seringueira, o convite para assumir o gerenciamento da SULANOR, nome recebido pela empresa a ser erguida em terras piauienses.

Quando ouvimos falar do Piauí, em 1983, como a maioria dos habitantes do centro-sul, não tínhamos a menor ideia de como era esse lugar.
Eu me lembrava apenas de uma referência a Teresina, a capital, lida num livro de infância, ``Viagens Através do Brasil'', onde o autor, Theobaldo Miranda Santos, creio, descrevia-a como \textit{``cidade-menina, pequenina e tímida, aninhada entre mangueiras''}. Sabia também que era a cidade mais quente do Brasil, se bem que eu não acreditasse que podia ser mais que Cuiabá, cidade onde vi um israelense desmaiar de calor, assim que deixou o ar refrigerado do hotel.

Logo que o convite foi apresentado ao Paulo, a mudança para o Piauí entrou na pauta das discussões em casa.
Aceitar, ou não aceitar? Paulo era o menos motivado.
Realmente, ele estava bem em Conceição do Araguaia.
O garimpo acabara, mas tínhamos uma poupança, ele tinha comprado um sítio, a assistência técnica às seringueiras e o viveiro de mudas nos sustentava e, além do mais, tínhamos pouquíssimos gastos, a não ser com as viagens de férias para São Paulo.
Entretanto, eu não estava segura de que nossas oportunidades de crescimento em Conceição fossem consistentes nem a curto e nem a médio prazo.


Essas cidades pequenas do interior do Brasil, parecem repetir um padrão no que toca às perspectivas de desenvolvimento.
Todo início de ano, começam a surgir boatos, aqui e ali, de que uma estrada vai ser aberta, ou que tal ou qual companhia vem se instalar nas imediações, ou que o governo tem um grande projeto para a região, e por aí vai até mais ou menos o meio do ano, quando as notícias começam a minguar, os comentários vão se dissipando, tudo recai no ramerrão de sempre e as esperanças são recolhidas para aguardar o inicio do próximo ano quando, esquecidas do desengano do ano anterior, as pessoas voltam de novo a se embalar com as notícias das grandes mudanças que irão acontecer neste ano que entra, sim, que agora é garantido, desse ano não passa, etc.,etc.,etc.
Era assim em Bonito, foi assim em Conceição e é assim em Cananéia, onde resido atualmente.
É verdade que todos esses lugares cresceram, hoje estão muito melhores do que quando os conheci.
Contudo, lá se foram vinte ou trinta anos e o crescimento que tiveram nesse período, à exceção de Bonito, foi, pode-se dizer, vegetativo, muito mais resultado do aumento demográfico do que de uma estratégia séria de incentivo ao desenvolvimento.
Eu achava que não tínhamos esse tempo todo para esperar.
Tomaz completara a quarta série do primário e ia para o primeiro grau maior.
Eu queria proporcionar-lhe um estudo de melhor qualidade.

Recorremos a todo mundo que conhecíamos para descobrir um pouco mais sobre Teresina, pois Miguel Alves, já sabíamos, era um lugarejo acanhado, muito pobre de recursos.
Até que um médico piauiense, morador há muitos em Conceição, Doutor Carlos, nos falou com entusiasmo da cidade e, principalmente, do Colégio São Francisco de Sales, um dos primeiros colégios jesuítas do Brasil, onde ele próprio estudara e que ele reputava ser um dos melhores da região nordeste.
Essa notícia me deixou feliz e até o famoso calor da cidade passou a ter importância secundária.


A essa altura, Paulo foi convidado a ir conhecer o lugar onde se instalaria o projeto e a verdade é que ele voltou também um pouco mais entusiasmado.
Teresina, ele contou, era uma cidade surpreendentemente grande, moderna, equivalente, em recursos, a qualquer das médias cidades paulistas.
Para inscrever as crianças no Colégio, precisaríamos submetê-las a uma avaliação prévia e apresentar algumas cartas de recomendação que atestassem os bons hábitos e a formação religiosa da família.

A ideia da mudança para o Piauí começou a tomar forma, vencendo o que ainda restava de resistência em Paulo.
Com algum pesar, mas reconhecendo que seria melhor para as crianças e quem sabe, para mim, aceitou a proposta da Sul América.

Pouco tempo depois, Paulo começou a viajar para Miguel Alves para tomar as primeiras providências necessárias à implantação do projeto.
Era um desafio na medida exata para empolgá-lo: erguer uma empresa de ponta num lugar onde as pessoas viviam do extrativismo e do escambo, literalmente.


O estado do Piauí, na época, era uma espécie de feudo mantido na pobreza e no atraso por uma pequena elite sobrevivente do tempo dos coronéis.
Nas pequenas cidades do interior, as pessoas viviam de trocar na feira o parco produto do seu trabalho: um pouco de mandioca, uma galinha, um cabrito, rapadura, uns litros de feijão de corda ou de farinha.
Mesmo nas áreas que viviam da quebra do coco babaçu, o pagamento era feito em espécie e era mísero.
Não corria dinheiro, a não ser aquele recebido pelos servidores da administração pública que, na sua maioria, não recebiam sequer o salário mínimo.
Aliás, o Estado era o maior e quase único empregador.
Em Teresina, que funcionava como um centro de abastecimento para algumas localidades do Pará, do Maranhão, do Mato Grosso e do Ceará e por isso ostentava um comércio tradicional, atacadista ou de ``estiva'' como eles diziam, ainda se podia encontrar serviço atrás dos balcões.
Porém as indústrias, que demandam mão de obra em grande escala, essas eram muito poucas.
A esperança, então, era conquistar um emprego público por meio de concurso, quando havia, ou, por indicação política.
Disso se beneficiavam os mandantes locais que, graças à atribuição de nomear e demitir, iam se perpetuando no poder, partilhando-o pelo histórico sistema de alternância.
Ora um, ora outro, a cada eleição cumpria-se a profecia conhecida de todos: ``a quatro anos de vacas gordas, seguem-se quatro de vacas magras'', ou seja, uma vez eleito, cada político vitorioso tratava de demitir os apaniguados do governo anterior, que ficavam à margem das benesses por amargos quatro anos, aguardando a oportunidade de reverter a situação.
As disputas eleitorais, portanto, eram acirradas e produziam tanto rupturas dramáticas entre aliados de primeira hora, quanto alianças inacreditáveis entre inimigos figadais de toda uma vida.
O importante era chegar ao poder.
Numa coisa, porém, os coronéis de todas as facções estavam de acordo: a grande fonte de onde brotava o manancial de recursos que sustentava todo o sistema ficava na Capital Federal e a palavra mágica para fazê-la jorrar era ``SECA''.
Ironicamente, a fazenda que a Sul América, a convite do governo do estado, ia estabelecer em Miguel Alves, demonstraria a falácia contida nos argumentos que por décadas justificaram toneladas de verbas enviadas a fundo perdido para o Piauí, para mitigar a fome do homem do campo vitimado pela seca.
Irrigado, o arrozal da SULANOR bateu recordes de produção.

Mas essa corrupta camarilha de pais da pátria nunca refletiu a verdadeira natureza do piauiense, como haveríamos de descobrir.
A brava gente do Piauí está entre os brasileiros, homens e mulheres, que mais migram nesse país atrás de trabalho.
Podemos encontrá-los nas mais distantes fronteiras, enfrentando os mais duros serviços e as mais difíceis condições de vida.
 Para os que lá preferem ficar, a necessidade de batalhar por um emprego no Estado resulta em grande apreço pelo estudo.
Essa seria outra auspiciosa descoberta para nós.

Em Janeiro de 1984, chegamos a Teresina.
As crianças tinham que se submeter ao teste de seleção no Colégio São Francisco de Sales.
As aulas começariam em Fevereiro.
Eu e eles estávamos ansiosos, na esperança de ver aprovado todo o nosso trabalho ao longo do tempo de Conceição.
No dia das tão esperadas provas, exceção feita à Paula que se assustou com a arquitetura pesada e as cores sombrias do interior do prédio centenário e desandou no choro, Tomaz, Marcelo e Fernando saíram-se muito bem, sem nenhuma dificuldade.
E mesmo a chorona, segundo a Coordenadora, apesar do medo, mostrou que dominava sem dificuldade o programa da sua série.
Aliviados, passamos ao problema seguinte: achar uma casa.

Reunindo o produto da venda do sítio de Conceição, mais a nossa poupança e uma ajuda que recebemos da Tia Maria, a irmã mais velha do meu sogro, nós podíamos pensar em comprar uma boa casa em Teresina e no melhor bairro da cidade, conhecido como o Jockey, por causa do clube de mesmo nome construído nessa área.
Eu estava morta de feliz com essa decisão do Paulo, porque ter uma casa para mim sempre foi prioridade.
Era uma questão de segurança e até de identidade.
Eu me sentia vulnerável e incompleta sem um lugar para morar que fosse nosso, que traduzisse o nosso jeito de ser e de viver e que não estivesse sujeito aos caprichos de um senhorio.
Paulo não conseguia entender essa necessidade e achava que, num início de vida, significava só um emprego improdutivo de capital.


Acabamos por achar uma casa grande e gostosa, de esquina, com um quintal enorme, sombreado por meia dúzia de frondosas mangueiras.
Esse era um dos encantos de Teresina: seus quintais e suas mangueiras.
E era também o modo como seus habitantes se defendiam do calor infernal.
Quando fechamos o negócio, mal podíamos acreditar que ela agora nos pertencia.
Como a antiga dona ainda demoraria algumas semanas para se mudar, tínhamos que nos contentar em contemplar nossa aquisição de longe para nos acostumarmos à ideia de que moraríamos naquela residência que, em Araraquara, pelo tamanho e pela aparência, seria uma casa de ricos.


\begin{figure}
\centering
\includegraphics[width=\linewidth]{34/bela-casa.png}
\caption{Nossa bela casa de Teresina.}
\end{figure}

\begin{figure}
\centering
\includegraphics[width=\linewidth]{34/brincando+cachorros.png}
\caption{ Brincando com os cachorros no quintal, pouco antes de ir para as aulas no Diocesano.}
\end{figure}

Em fevereiro, as aulas começaram.
Enquanto as crianças se adaptavam à nova escola, eu tratava de reorganizar a rotina da família e conhecer a cidade.
Uma manhã, quando tirava o carro da garagem, tive a atenção chamada por alguém que, em altos brados, tentava orientar o motorista meio perdido de um caminhão de entregas.
Era uma vizinha, muito despachada, que logo tratou de partilhar o contratempo comigo.
Há horas, disse, esperava pelo tal caminhão que dava voltas pelo bairro sem atinar com o endereço dela.
Enquanto o motorista manobrava para embicar na nossa ruazinha, ela se apresentou: Maria Luzia.
Ela e o marido, Chico, haveriam de se tornar os nossos maiores amigos na cidade.
Amigos, compadres, quase irmãos.


Adiante da casa de Maria Luzia, moravam Lúcia e Castelo.
Ele, um médico baixinho, gorducho, simpático e bonachão.
Ela, ao contrário, me pôs de sobreaviso ao se fazer logo muito chegada, muito íntima, muito curiosa sobre o trabalho do Paulo e nos convidando para isso e aquilo, sem admitir recusa.
O tempo confirmaria que a tal mulherinha precisava ser mantida a uma distância estratégica e com luvas de pelica.
Era maníaca, depressiva e encrenqueira.
Mas, o pobre marido merecia toda nossa consideração.

A adaptação das crianças estava sendo rápida.
A casa já funcionava dentro de uma rotina normal, Iracema e Deusa tinham chegado de Conceição e tudo caminhava da melhor maneira, quando um dia Fernando começou a queixar-se de uma dor forte num lado do abdômen.
Médico, só conhecíamos o Castelo.
Por intermédio de Lúcia, conseguimos localizá-lo e, sem demora, ele veio até em casa.
Bastou um rápido exame para que ele chegasse ao diagnóstico:

\textit{``-- É um apêndice agudo e Fernando tem que ser operado.
Vou encaminhá-los a um colega, cirurgião pediátrico.''}

Internamos Fernando à noitinha no hospital que, segundo o cirurgião, era o único dotado de um centro cirúrgico seguro.
Nervoso, Paulo já explodiu na entrada quando quiseram deitar nosso menino numa maca cujo lençol exibia manchas de sangue.
Ao chegar ao quarto e abrir a porta do frigobar, fui surpreendida por um punhado de baratas vivas e estonteadas pela luz que subitamente invadiu o seu doce refrigério.
Eu queria chorar de pânico, imaginando que tipo de segurança ofereceria o centro cirúrgico de um hospital daqueles.
Mas tive que me conter porque não havia o que fazer, não havia tempo para pensar em remoção para São Paulo e porque se eu desmoronasse só Deus sabe qual seria a reação do Paulo.
Só restava rezar, rezar, rezar.

Quatro horas da manhã, um crioulo magro e retinto entrou quarto adentro empurrando a maca que levaria Fernando até a sala de cirurgia.
Acompanhamos nosso filho já meio sedado, nós próprios anestesiados pela noite em claro e pelo medo que nos empastava a consciência.
Vimos a maca desaparecer por uma portinhola baixa, no fim de um corredor, vedada por uma cortina de tiras de plástico.
Sentamo-nos num banco encostado à parede para esperar.
 Éramos eu, Paulo e o crioulo que agora, munido de rodinho, pano e balde, começava a esfregar o piso de cerâmica vermelha.
Ficamos acompanhando seus lentos movimentos, hipnotizados, mudos, rendidos.
Eu tentava não pensar.
De repente, uma campainha soou.
O crioulo encostou o rodinho e foi até a portinhola onde ficou aguardando.
Logo, alguém empurrou uma maca vazia através da cortina que ele, no seu passo arrastado, veio trazendo pelo corredor até desaparecer, lá no fim.
Minutos depois ressurgiu, transportando sobre a maca uma grávida chorosa, contorcendo-se em dores.
Levou-a até a tal portinhola, empurrou-a através das tiras de plástico e, imperturbável, retomou seu trabalho.
Metodicamente, mergulhou o pano sujo no balde, chacoalhou-o e, enrolando-o outra vez no rodinho, continuou a desenhar largas faixas de um castanho claro sobre as lajotas vermelhas.
Uma hora depois, a campainha soou novamente e a cena se repetiu.
O crioulo largou o rodinho num canto e caminhou até a portinhola.
Quando se virou para voltar, trazia nos braços um embrulhinho de flanela, dentro do qual adivinhei um bebê.
Quando passou por nós, um vagido de gatinho novo confirmou meu acerto.
Aquilo foi demais.
O homem sequer lavava a mão entre uma atribuição e outra.
Eu me desesperei, queria saber do meu filho, estava prestes a ter um ataque quando, no final do corredor, ao invés do crioulo, o que apareceu foi o rosto redondo e sorridente do Castelo.
Seguido da Lúcia, da Maria Luzia, da mãe da Maria Luzia, da Dona Terezinha, a senhora de quem tínhamos comprado a casa e mais um monte de pessoas que eu mal pude divisar através das lágrimas que me subiram aos olhos.
Foi precisamente aí que eu aprendi a querer bem aquela terra e aquela gente.
Castelo embarafustou pelo centro cirúrgico e em poucos minutos voltou acenando com um vidrinho cheio de álcool onde boiava o apêndice extirpado do Fernando.

\textit{``-- A operação transcorreu muito bem''}, garantiu-nos ele, \textit{``e logo o Fernando vai ser levado ao quarto.
Está na sala de recuperação, acordando da anestesia''}.

Nem me importei quando o indefectível crioulo apareceu empurrando o Fernando na maca.
Estava aliviada, feliz.
Meu filho vinha bravo, dolorido, mas são e salvo.
Era tudo o que importava.


